\section{Bounded primitives for Hochschild cocycles}\label{sec:cohomology}

We first apply Theorem~\ref{thm:liouville} to separately normal
cocycles on a type $\mathrm{II}_1$ algebra with separable predual and a
faithful normal tracial state.  The resulting uniform bound for a
primitive will allow us first to remove separability and then to
combine central summands.  Throughout this
section, normality of a linear map means ultraweak continuity, and
separate normality has the corresponding meaning for multilinear maps.

\subsection{Normality and continuity on bounded balls}

The continuity needed by Theorem~\ref{thm:liouville} is the tracial
form of Akemann's continuity theorem
\cite[Corollary~III.10]{Akemann1967}. We include a direct proof that
allows a center and does not require a separable predual.
The stronger joint multilinear continuity result for type
$\mathrm{II}_1$ factors is \cite[Theorem~4.4]{CPSS2003}.

\begin{lemma}\label{lem:normal-two-continuity}
Let $M$ be a finite von Neumann algebra with faithful normal tracial
state $\tau$, and let $B:M\to M$ be bounded and normal.  For every
$R<\infty$ and $\varepsilon>0$ there is $\delta>0$ such that
\[
 \norm{x}\le R,\quad \norm{x}_2<\delta
 \quad\Longrightarrow\quad \norm{B(x)}_2<\varepsilon.
\]
Consequently $B$ is $2$-norm continuous on every operator norm ball.
\end{lemma}

\begin{proof}
We argue by contradiction.  A failure of continuity will give elements
with mutually orthogonal two-sided supports whose images stay uniformly
away from zero in $2$-norm.  Random signs will then contradict the
boundedness of $B$.

If the first assertion fails, there are $R>0$, $\eta>0$, and a sequence
$(x_j)$ such that
\[
 \norm{x_j}\le R,\qquad \norm{x_j}_2\longrightarrow0,
 \qquad \norm{B(x_j)}_2\ge\eta.
\]
Choose $\alpha>0$ with $\norm{B}\alpha<\eta/4$.  In the polar
decomposition $x_j=v_j|x_j|$, set
\[
 r_j=1_{(\alpha,\infty)}(|x_j|),\qquad
 y_j=v_j|x_j|r_j,\qquad
 \ell_j=v_jr_jv_j^*,\qquad e_j=r_j\vee\ell_j.
\]
Here $r_j$ and $\ell_j$ are the right and left support projections of
$y_j$.  In particular,
\[
 y_j=e_jy_je_j,\quad \norm{y_j}\le R,\quad
 \norm{x_j-y_j}\le\alpha,\quad
 \norm{B(y_j)}_2\ge 3\eta/4.
\]
The trace and the spectral truncation give
\[
 \tau(e_j)\le\tau(r_j)+\tau(\ell_j)
 =2\tau(r_j)\le2\alpha^{-2}\norm{x_j}_2^2.
\]
Pass to a subsequence for which $\sum_j\tau(e_j)<\infty$, and put
$p_j=\bigvee_{h\ge j}e_h$.  Then $p_j$ decreases to zero: indeed,
$\tau(p_j)\le\sum_{h\ge j}\tau(e_h)\to0$, and $\tau$ is faithful
and normal.

For fixed $j$ and $h>j$, the compressions
\[
 (p_j-p_h)y_j(p_j-p_h)
\]
converge ultraweakly to $y_j$ as $h\to\infty$.  Applying $B$ preserves
this convergence.  The $2$-norm is ultraweakly lower semicontinuous,
as is seen from
\[
 \norm{y}_2=
 \sup\{ |\tau(b^*y)|:b\in M,\ \norm{b}_2\le1\}.
\]
Starting with $j_1=1$, we may therefore choose successively
$j_{s+1}>j_s$ so that
\[
 t_s=(p_{j_s}-p_{j_{s+1}})y_{j_s}(p_{j_s}-p_{j_{s+1}}),
 \qquad \norm{B(t_s)}_2\ge\eta/2.
\]
The $t_s$ have pairwise orthogonal two-sided supports and
$\norm{t_s}\le R$.  Hence
$\norm{\sum_{s=1}^N\epsilon_st_s}\le R$ for every choice of signs
$\epsilon_s\in\{-1,1\}$.  Taking independent uniform signs, Hilbert
space orthogonality of these scalar random variables yields
\[
 \frac{N\eta^2}{4}
 \le\sum_{s=1}^N\norm{B(t_s)}_2^2
 =\E\norm{B\Bigl(\sum_{s=1}^N\epsilon_st_s\Bigr)}_2^2
 \le\norm{B}^2R^2,
\]
which is impossible for large $N$.  This proves the asserted modulus.
Applying it to differences, which lie in the ball of radius $2R$,
gives continuity on the ball of radius $R$.
\end{proof}

\subsection{Averaging the last input}

\begin{proposition}\label{prop:separable-primitive}
Let $M$ be of type $\mathrm{II}_1$ with separable predual and faithful
normal tracial state $\tau$.  If $k\ge2$ and
$f\in C_b^k(M,M)$ is a separately normal cocycle, then there is
$g\in C_b^{k-1}(M,M)$ such that
\begin{equation}\label{eq:separable-primitive-bound}
 dg=f,\qquad \norm{g}\le\norm{f}.
\end{equation}
\end{proposition}

\begin{proof}
We first construct an averaged cochain that is a right module map in
its last input.  Evaluation at $1$ will give a $(k-1)$-cochain, and averaging
the original equation $df=0$ will identify its differential.

Choose the symmetric countably supported measure $\mu$ for which
Theorem~\ref{thm:liouville} applies, and write
$\mu=\sum_{s\ge1}p_s\delta_{u_s}$.  For a bounded $k$-cochain
$\Phi$, define
\begin{equation}\label{eq:last-input-average}
 (P\Phi)(\mathbf a,x)
 =\sum_{s\ge1}p_s\Phi(\mathbf a,xu_s)u_s^*,
 \qquad \mathbf a=(a_1,\ldots,a_{k-1}).
\end{equation}
The sum converges in operator norm, and $P$ is a contraction for the
cochain norm.  If $W_j=U_1\cdots U_j$, where the $U_i$ are independent
with law $\mu$, and $W_0=1$, induction gives
\[
 (P^j f)(\mathbf a,x)
 =\E\bigl[f(\mathbf a,xW_j)W_j^*\bigr].
\]
Set
\[
 F_n=\frac1n\sum_{j=0}^{n-1}P^j f.
\]

We use compactness for pointwise ultraweak convergence.  For each fixed
tuple $\mathbf x=(x_1,\ldots,x_k)$, all values $F_n(\mathbf x)$ lie
in the ultraweakly compact ball of radius
$\norm{f}\prod_i\norm{x_i}$.  The product of these balls, indexed by
all tuples in $M^k$, is compact.  Multilinearity is closed for this
topology, since addition and scalar multiplication are ultraweakly
continuous.  Thus some subnet of $(F_n)$ converges at every tuple to
a bounded multilinear map $F$ with $\norm{F}\le\norm{f}$.

On any common cochain norm ball, $P$ is continuous for this topology.
To see this, suppose $\Phi_\lambda\to\Phi$ pointwise ultraweakly and
$\norm{\Phi_\lambda},\norm{\Phi}\le C$.  For a fixed tuple
$(\mathbf a,x)$ and a normal functional $\rho\in M_*$, truncate
\eqref{eq:last-input-average} after $L$ terms.  The finite sum passes
to the limit, and the remaining difference has absolute value at most
\[
 2C\norm{\rho}\norm{x}
 \prod_{i=1}^{k-1}\norm{a_i}\sum_{s>L}p_s.
\]
Letting $L\to\infty$ proves the claim.  Since
\[
 PF_n-F_n=\frac{P^nf-f}{n},\qquad
 \norm{PF_n-F_n}\le\frac{2\norm{f}}n,
\]
the subnet limit satisfies $PF=F$.

Fix $\mathbf a$, and put $B(x)=f(\mathbf a,x)$.  This is a bounded
normal linear map.  For any $R<\infty$ and $\varepsilon>0$,
Lemma~\ref{lem:normal-two-continuity} provides $\delta>0$ such that
\[
 \norm{x}\le R,\quad \norm{x}_2<\delta
 \quad\Longrightarrow\quad
 \norm{B(xv)v^*}_2<\varepsilon
 \quad\hbox{for every }v\in\U(M).
\]
Indeed, right multiplication by a unitary preserves both the operator
norm and the $2$-norm.  The same bound holds after averaging over any
$W_j$ and after taking the Ces\`aro average.  Ultraweak lower
semicontinuity of the $2$-norm then gives
$\norm{F(\mathbf a,x)}_2\le\varepsilon$ under these input bounds.
Using $\varepsilon/2$ in place of $\varepsilon$ proves the required
continuity modulus for $T(x)=F(\mathbf a,x)$.  The equality $PF=F$
is exactly the harmonicity identity for this map.  Consequently
Theorem~\ref{thm:liouville} gives
\begin{equation}\label{eq:averaged-right-module}
 F(\mathbf a,x)=h(\mathbf a)x,
 \qquad h(\mathbf a)=F(\mathbf a,1).
\end{equation}
In particular $h$ is a bounded $(k-1)$-cochain and
$\norm{h}\le\norm{f}$.

It remains to obtain its coboundary identity.  Let
$\nu_n=n^{-1}\sum_{j=0}^{n-1}\mu^{*j}$ and
$h_n(\mathbf a)=F_n(\mathbf a,1)$.  Evaluate the original equation
$df=0$ at $(a_1,\ldots,a_k,v)$, multiply the output on the right by
$v^*$, and average with respect to $v\sim\nu_n$.  This gives the
exact equality
\begin{align}
 0={}&a_1h_n(a_2,\ldots,a_k)
 +\sum_{j=1}^{k-1}(-1)^j
 h_n(a_1,\ldots,a_ja_{j+1},\ldots,a_k)\notag\\
 &+(-1)^k F_n(a_1,\ldots,a_{k-1},a_k)
 +(-1)^{k+1}f(a_1,\ldots,a_k).
 \label{eq:original-cocycle-average}
\end{align}
The final term is unchanged because $vv^*=1$.  Along the chosen
subnet, each displayed evaluation converges ultraweakly; multiplication
by a fixed $a_i$ preserves that convergence.  Passing to the limit
and using \eqref{eq:averaged-right-module} therefore yields
\[
 0=(dh)(a_1,\ldots,a_k)+(-1)^{k+1}f(a_1,\ldots,a_k).
\]
Thus $g=(-1)^kh$ satisfies \eqref{eq:separable-primitive-bound}.
This calculation averages $df$ at the displayed tuples; it requires
neither a commutation relation between $P$ and $d$ nor separate
normality of the limiting cochain $F$.
\end{proof}

\subsection{Removing the separability assumption}

Uniform primitive bounds and pointwise ultraweak compactness were
used to remove separability in \cite[Section~7]{CPSS2003}.
As in that argument, conditional expectations allow the approximating subalgebras
to be chosen independently for each finite set of inputs.

\begin{proposition}\label{prop:tracial-primitive}
Let $M$ be any type $\mathrm{II}_1$ von Neumann algebra admitting a
faithful normal tracial state.  Every separately normal cocycle
$f\in C_b^k(M,M)$, $k\ge2$, has a bounded primitive $g$ with
$\norm{g}\le\norm{f}$.
\end{proposition}

\begin{proof}
Order the finite subsets $E\subset M$ by inclusion.  For each $E$,
choose a countably generated unital von Neumann subalgebra $N_E$
containing $E$, the finite set $f(E^k)$, and unital matrix systems
of every size $m\ge2$.  Such systems exist in a type $\mathrm{II}_1$
algebra: the unit can be divided into $m$ equivalent orthogonal
projections~\cite[Examples~III.2.5.4(iii)]{Blackadar2006}, and partial isometries between them supply the matrix
units.  Adjoining one system for each $m$ adds only countably many
generators.

The restriction of the trace to $N_E$ is faithful and normal.  The
tracial Hilbert space $L^2(N_E)$ is separable: $*$-polynomials with
coefficients in $\mathbb Q+i\mathbb Q$ in a countable set of generators are dense in the
generated $C^*$-algebra in norm, and bounded strong approximation
makes that algebra $2$-norm dense in $N_E$.  The faithful tracial
representation thus gives $N_E$ separable predual.  Moreover $N_E$
has no type I central part.  A nonzero homogeneous type $\mathrm{I}_n$
central summand would inherit a unital copy of $M_m(\C)$ for every
$m$, whereas its $n$-dimensional fibers cannot admit a unital copy
when $m>n$.  Every finite type I algebra decomposes into its
homogeneous type $\mathrm{I}_n$ summands.  Since $N_E$ is finite,
it is therefore of type $\mathrm{II}_1$.

Let $\mathcal E_E:M\to N_E$ be the normal trace-preserving
conditional expectation. Its existence and bimodularity are classical;
see \cite[Th\'eor\`eme~8]{Dixmier1953} and
\cite[Theorem~1]{Umegaki1954}. Define
\[
 f_E=\mathcal E_E\circ f\big|_{N_E^k}.
\]
This map is separately normal and has norm at most $\norm{f}$.
For inputs in $N_E$, bimodularity of $\mathcal E_E$ gives
$df_E=\mathcal E_E(df)=0$.  By
Proposition~\ref{prop:separable-primitive}, choose
$g_E\in C_b^{k-1}(N_E,N_E)$ such that
$dg_E=f_E$ and $\norm{g_E}\le\norm{f}$.  Extend it to a cochain on
$M$ by
\[
 G_E(x_1,\ldots,x_{k-1})
 =g_E(\mathcal E_Ex_1,\ldots,\mathcal E_Ex_{k-1}).
\]
This extension has the same norm bound.  If
$(a_1,\ldots,a_k)\in E^k$, all its inputs and adjacent products
belong to $N_E$.  Expanding the differential consequently gives
\begin{equation}\label{eq:finite-set-exactness}
 (dG_E)(a_1,\ldots,a_k)
 =(dg_E)(a_1,\ldots,a_k)
 =\mathcal E_Ef(a_1,\ldots,a_k)
 =f(a_1,\ldots,a_k),
\end{equation}
where the final equality uses $f(E^k)\subset N_E$.

The common norm bound and the same product compactness used above
give a pointwise ultraweakly convergent subnet of $(G_E)$, with limit
$g\in C_b^{k-1}(M,M)$ and $\norm{g}\le\norm{f}$.  For a fixed
$k$-tuple $(a_1,\ldots,a_k)$, the indices of this subnet eventually
contain $E_0=\{a_1,\ldots,a_k\}$.  Equation~\eqref{eq:finite-set-exactness}
therefore holds eventually along the subnet.  Passing to the limit
in its finitely many coboundary terms gives $dg=f$ at this tuple.
This holds at every tuple, proving the proposition.  The choices of
$N_E$ and $\mathcal E_E$ can be made independently for different
$E$; eventual exactness is supplied by the directed finite subsets.
\end{proof}

\subsection{Central cuts and mixed inputs}

Uniform bounds and central decomposition also enter the tensor-product
argument of \cite[Section~4]{PopSmith2010}. In the present setting we
use an explicit homotopy to correct all inputs outside a given central
summand before assembling the primitives.

For a central projection $z\in M$, regard $zM$ as an $M$-bimodule
under multiplication and put $q=1-z$.  In degree zero we use
$C_b^0(M,zM)=zM$ and $(d\xi)(a)=a\xi-\xi a$.
On cochains with values in
$zM$, define the central cut by
\[
 (\operatorname{cut}_z\psi)(a_1,\ldots,a_n)
 =\psi(za_1,\ldots,za_n).
\]
In degree zero this means $\operatorname{cut}_z=\id$ on $zM$.
Centrality and the fact that $q$ annihilates this module show that
$\operatorname{cut}_z$ commutes with $d$.
For a cocycle $\psi$ with values in $zM$, a primitive for its
restriction to $zM$, pulled back to $M$, accounts only for
$\operatorname{cut}_z\psi$.
The following homotopy supplies the remaining difference
$\psi-\operatorname{cut}_z\psi$ on arbitrary inputs, with a bound
independent of $z$.

\begin{lemma}\label{lem:central-homotopy}
For $n\ge2$ and $\psi\in C_b^n(M,zM)$, set
\begin{equation}\label{eq:central-homotopy-formula}
 \begin{split}
 (J_z\psi)(a_1,\ldots,a_{n-1})
 =\sum_{i=1}^{n-1}(-1)^i\psi(&za_1,\ldots,za_{i-1},
 q,qa_i,\\[-2pt]
 &a_{i+1},\ldots,a_{n-1}).
 \end{split}
\end{equation}
Set $J_z=0$ in degrees zero and one.  Then, in every degree,
\begin{equation}\label{eq:central-homotopy-identity}
 dJ_z+J_zd=\id-\operatorname{cut}_z.
\end{equation}
For $n\ge1$ the degree-$n$ operator satisfies
\begin{equation}\label{eq:central-homotopy-bound}
 \norm{J_z\psi}\le(n-1)\norm{\psi}.
\end{equation}
\end{lemma}

\begin{proof}
The norm estimate follows directly from
\eqref{eq:central-homotopy-formula}.  To verify the identity, expand
each input as $a_i=za_i+qa_i$.  By multilinearity it suffices to
use tuples whose entries each belong to $zM$ or to $qM$.  On such
a tuple with first $qM$ entry at position $r$, formula
\eqref{eq:central-homotopy-formula} consists of the single term
inserting $q$ immediately before that entry, with sign $(-1)^r$;
on an all-$zM$ tuple it is zero.

Consider such an $n$-tuple with first $qM$ entry $a_r$.
In $(J_zd)\psi$, merging the inserted $q$ with $a_r$ has total
sign $(-1)^r(-1)^r=1$ and contributes
$\psi(a_1,\ldots,a_n)$.  A merger strictly before the insertion
has the same arguments as its counterpart in $(dJ_z)\psi$,
but the latter inserts at position $r-1$, so their signs are
opposite.  A merger strictly after the pair $(q,a_r)$ also has
the same arguments as its counterpart in $(dJ_z)\psi$, and its
merger index is larger by one, again reversing the sign.  The
merger of $a_{r-1}$ with the inserted $q$, when $r>1$, vanishes
because $zMqM=0$; the corresponding merger across $a_{r-1},a_r$
in $(dJ_z)\psi$ vanishes for the same reason.

The left outer terms cancel when $r>1$, with insertion indices
$r$ and $r-1$.  When $r=1$, both left outer terms vanish because
$qM$ annihilates the output module.  For the right outer terms,
if $r<n$ the same insertion index and the two consecutive
differential degrees give opposite signs.  If $r=n$, the last
input is in $qM$, so both right outer terms vanish.  Thus the
only surviving term is $\psi(a_1,\ldots,a_n)$, as required
because the cut of this tuple is zero.  On an all-$zM$ tuple,
both sides of \eqref{eq:central-homotopy-identity} vanish.

For completeness, in degree one the formula reads
\[
 (J_zd\psi)(a)=-(d\psi)(q,qa)=\psi(qa)
 =\psi(a)-\psi(za),
\]
since both outer module actions in $(d\psi)(q,qa)$ vanish.
In degree zero both sides are zero by the definitions.  This
proves the identity in all degrees.
\end{proof}

In particular, if $\psi\in C_b^k(M,zM)$ is a cocycle, then
\begin{equation}\label{eq:central-cocycle-correction}
 d(J_z\psi)=\psi-\operatorname{cut}_z\psi,
 \qquad \norm{J_z\psi}\le(k-1)\norm{\psi}.
\end{equation}
This bound is independent of the algebra and the central projection.

\subsection{Completion of the proof}

\begin{proof}[Proof of Theorem~\ref{thm:main}]
The zero algebra is immediate, so assume $M\ne0$.
Fix $k\ge2$ and an arbitrary bounded cocycle
$f_{\mathrm{orig}}\in C_b^k(M,M)$.  By the normal reduction in
Proposition~\ref{prop:classical}, there are a separately normal
cocycle $f$ and a bounded $(k-1)$-cochain $b$ such that
\[
 f_{\mathrm{orig}}=f+db.
\]
Let $z_{\mathrm{II}_1}$ be the central projection of the type
$\mathrm{II}_1$ summand, and let $z_0=1-z_{\mathrm{II}_1}$.
We first partition $z_{\mathrm{II}_1}$ into central projections
whose algebras admit faithful normal tracial states.

In the abelian algebra $Z(z_{\mathrm{II}_1}M)$, choose a maximal
orthogonal family $(z_i)_{i\in I}$ of nonzero supports of normal
states.  Its join is $z_{\mathrm{II}_1}$: a nonzero complementary
projection would support a further normal state, contradicting
maximality.  A state is faithful on its support, so each
$z_iZ(M)$ has a faithful normal state $\varphi_i$.  Composing
$\varphi_i$ with the faithful normal center-valued trace of $z_iM$
gives a faithful normal tracial state $\tau_i$ on $z_iM$;
see \cite[Theorem~III.2.5.7 and Corollary~III.2.5.8]{Blackadar2006}.
The family $I$ may be uncountable.  Omit $z_0$ if it is zero,
and otherwise keep the entire complementary summand as one piece.

For each piece define the restricted cocycle
\[
 f_i(x_1,\ldots,x_k)=z_i f(x_1,\ldots,x_k),
 \qquad x_j\in z_iM.
\]
Multiplication by a central projection is a bimodule map, so
$f_i$ is a separately normal cocycle and
$\norm{f_i}\le\norm{f}$.  For $i\in I$,
Proposition~\ref{prop:tracial-primitive} supplies a primitive
$b_i\in C_b^{k-1}(z_iM,z_iM)$ with
\[
 db_i=f_i,\qquad \norm{b_i}\le\norm{f}.
\]
If $z_0\ne0$, the vanishing result for algebras without a type
$\mathrm{II}_1$ central summand in
Proposition~\ref{prop:classical} supplies a bounded primitive
$b_0$ of $f_0$.  Its norm is some finite number.

Pull each primitive back to $M$ through the homomorphism
$x\mapsto z_ix$:
\[
 \widetilde b_i(x_1,\ldots,x_{k-1})
 =b_i(z_ix_1,\ldots,z_ix_{k-1}).
\]
The outer module actions also pull back, since the outputs lie
in $z_iM$.  Thus
\[
 d\widetilde b_i(x_1,\ldots,x_k)
 =z_i f(z_ix_1,\ldots,z_ix_k)
 =\operatorname{cut}_{z_i}(z_if)(x_1,\ldots,x_k).
\]
Define cochains with values in $z_iM$ on all of $M$ by
\[
 H_i=\widetilde b_i+J_{z_i}(z_if).
\]
Equation~\eqref{eq:central-cocycle-correction} gives
\[
 dH_i=z_if,
 \qquad
 \norm{H_i}\le k\norm{f}\quad(i\in I),
\]
and, when present,
\[
 \norm{H_0}\le\norm{b_0}+(k-1)\norm{f}.
\]
There is consequently a single finite constant $C$ bounding the
norms of all these cochains; one can take the maximum of the two
displayed bounds, omitting the second when $z_0=0$.

For an arbitrary central partition of unity, the map
\[
 M\longrightarrow\prod_i^{\ell^\infty}z_iM,
 \qquad x\longmapsto(z_ix)_i,
\]
is an isometric $*$-isomorphism onto the bounded direct product.
For surjectivity, the finite partial sums of a bounded family
$(x_i)_i$, $x_i\in z_iM$, form a net directed by inclusion of the
finite index sets.  This net converges strongly, since the central
supports are orthogonal; its limit belongs to $M$.
For each tuple $(x_1,\ldots,x_{k-1})$, the family
$(H_i(x_1,\ldots,x_{k-1}))_i$ belongs to this product and has norm
at most $C\prod_j\norm{x_j}$.  Let $G(x_1,\ldots,x_{k-1})\in M$
be the corresponding element.  Coordinatewise multilinearity
shows that $G$ is a bounded multilinear cochain with
$\norm{G}\le C$.  Central projection commutes with every term
of its coboundary, so
\[
 z_i(dG)=dH_i=z_if
 \qquad\hbox{for every central piece }i.
\]
Since the projections have join $1$, it follows that $dG=f$.
Finally,
\[
 d(G+b)=f+db=f_{\mathrm{orig}}.
\]
We have obtained an actual bounded $(k-1)$-cochain whose
coboundary is the given cocycle.  Hence
$H_b^k(M,M)=0$ for every complex von Neumann algebra $M$ and
every $k\ge2$, with the actual image of $d$ in the definition of
cohomology and with no complete-boundedness hypothesis.
\end{proof}
